\documentclass[12pt]{article}
\usepackage[utf8]{inputenc}
\usepackage[T1]{fontenc}
\usepackage[margin=1in]{geometry}
\usepackage{setspace}
\usepackage{xcolor}
\usepackage[colorlinks=true, urlcolor=blue!60!black, linkcolor=blue!60!black]{hyperref}
\setlength{\parindent}{0em}
\setlength{\parskip}{0.6em}
\renewcommand{\baselinestretch}{1.0}
\begin{document}
\begin{center}
{\LARGE \textit{Lujan v. Defenders of Wildlife}, 504 U.S. 555 (1992)\par}
\vspace{0.5em}
{\large Case Brief --- Two-Page Edition\par}
\end{center}

\section{Memory Jogger}
The three-part standing test: vague ``some day'' travel plans are not imminent injury, and Congress cannot create standing for generalized procedural grievances (\href{https://supreme.justia.com/cases/federal/us/504/555/}{Justia}).

\section{Detailed Case Facts}
A 1986 regulation limited the Endangered Species Act's consultation duty to actions in the United States or on the high seas. Defenders of Wildlife challenged it, relying on members who had once visited Egypt and Sri Lanka near U.S.-funded projects and intended to return ``in the future.''

\section{Procedural History}
The Eighth Circuit found standing and held that the consultation duty applies abroad.

\section{Judicial Votes}
Reversed 7--2 (6--3 on standing). Scalia wrote for the Court (Part III-B for a plurality); Kennedy, joined by Souter, concurred in part; Stevens concurred in the judgment; Blackmun, joined by O'Connor, dissented.

\section{Holding}
Standing requires a concrete, particularized, actual or imminent injury, fairly traceable to the defendant and likely redressable. Congress cannot confer standing on ``any person'' to enforce an abstract procedural duty.

\section{Analysis of Opinions}
\textbf{Majority:} past visits and ``some day'' intentions fail; nexus theories would allow standing based on ``pure speculation and fantasy''; and letting citizens sue over executive compliance would transfer the President's Take Care duty to courts. \textbf{Kennedy:} Congress may define new injuries if it ties them to a class of plaintiffs. \textbf{Stevens:} standing existed, but the statute does not apply abroad. \textbf{Blackmun:} demanding plane tickets is ``an empty formality.''

\section{Examples: Future Applications}
\textbf{Same-Side:} a photographer who merely hopes to go back to Brazil ``someday'' has no standing. \textbf{Opposite-Side:} a researcher booked to study the very reef to be dredged has it. \textbf{Fence-Sitter:} a consumer with an undisseminated credit-report error suing under a statutory right, the question later taken up in \textit{Spokeo} and \textit{TransUnion}.

\section{Critique}
Historians note that ``injury in fact'' is a modern invention and that early Congresses authorized informer suits; Sunstein and Pierce argue the decision limits Congress without textual warrant.

\section{Key Quotations}
\begin{quote}
``Over the years, our cases have established that the irreducible constitutional minimum of standing contains three elements. First, the plaintiff must have suffered an `injury in fact'---an invasion of a legally protected interest which is (a) concrete and particularized, ... and (b) `actual or imminent, not ``conjectural'' or ``hypothetical,'' ' ... Second, there must be a causal connection between the injury and the conduct complained of---the injury has to be `fairly ... trace[able] to the challenged action of the defendant, and not ... th[e] result [of] the independent action of some third party not before the court.' ... Third, it must be `likely,' as opposed to merely `speculative,' that the injury will be `redressed by a favorable decision.' ''
\end{quote}
\noindent --- Justice Scalia, for the Court, 504 U.S. at 560--61.

\end{document}
